\par\Needspace{6\baselineskip}\medskip
\hypertarget{native-089-article070}{}\pdfbookmark[2]{70.}{article-070}
\begin{center}70.\end{center}
Ne quid eorum, quae ad calculum locorum geocentricorum requiruntur, hic desideretur, quaedam adhuc de \emph{parallaxi} atque \emph{aberratione} adiicienda sunt. Methodum quidem supra iam descripsimus, secundum quam locus parallaxi affectus, i. e. cuilibet in superficie terrae puncto respondens, immediate maximaque facilitate determinari potest: sed quum in methodo vulgari in art. 62 et sequ. tradita locus geocentricus ad terrae centrum referri soleat, in quo casu a parallaxi liber dicitur, methodum peculiarem pro determinanda parallaxi, quae est inter utrumque locum differentia, adiicere oportebit.

\SourceLeaf{090}Sint corporis coelestis longitudo et latitudo geocentrica respectu centri terrae $\lambda$, $\beta$; eaedem respectu puncti cuiusvis in superficie terrae $l$, $b$; distantia corporis a terrae centro $r$, a puncto superficiei $\Delta$; denique respondeat in sphaera coelesti ipsi zenith huius puncti longitudo $L$, latitudo $B$, designeturque radius terrae per $R$. Sponte iam patet, omnes aequationes art. 62 etiam hic locum esse habituras; sed notabiliter contrahi poterunt, quum $R$ hic exprimat quantitatem prae $r$ et $\Delta$ tantum non evanescentem. Ceterum eaedem aequationes manifesto etiamnum valebunt, si $\lambda$, $l$, $L$ pro longitudinibus ascensiones rectas, atque $\beta$, $b$, $B$ pro latitudinibus declinationes exprimunt. In hoc casu $l-\lambda$, $b-\beta$ erunt parallaxes ascensionis rectae et declinationis, in illo vero parallaxes longitudinis et latitudinis. Quodsi iam $R$ ut quantitas primi ordinis tractatur, eiusdem ordinis erunt $l-\lambda$, $b-\beta$, $\Delta-r$, neglectisque ordinibus superioribus e formulis art. 62 facile derivabitur:
\[
\begin{aligned}
\mathrm{I.}\quad l-\lambda&=\frac{R\cos B\sin(\lambda-L)}{r\cos\beta}\\[.5em]
\mathrm{II.}\quad b-\beta&=\frac{R\cos B\cos\beta}{r}\{\operatorname{tang}\beta\cos(\lambda-L)-\operatorname{tang}B\}\\[.5em]
\mathrm{III.}\quad \Delta-r&=-R\cos B\sin\beta\{\operatorname{cotang}\beta\cos(\lambda-L)+\operatorname{tang}B\}.
\end{aligned}
\]
\noindent Accipiendo angulum auxiliarem $\theta$ ita ut fiat $\operatorname{tang}\theta=\dfrac{\operatorname{tang}B}{\cos(\lambda-L)}$, aequationes II, III formam sequentem nanciscuntur:
\[
\begin{aligned}
\mathrm{II.}\quad b-\beta&=\frac{R\cos B\cos(\lambda-L)\sin(\beta-\theta)}{r\cos\theta}=\frac{R\sin B\sin(\beta-\theta)}{r\sin\theta}\\[.5em]
\mathrm{III.}\quad \Delta-r&=-\frac{R\cos B\cos(\lambda-L)\cos(\beta-\theta)}{\cos\theta}=-\frac{R\sin B\cos(\beta-\theta)}{\sin\theta}\,\cdot
\end{aligned}
\]
\noindent Ceterum patet, ut in I et II $l-\lambda$ et $b-\beta$ in minutis secundis obtineantur, pro $R$ accipi debere parallaxem mediam solarem in minutis secundis expressam; in III vero pro $R$ eadem parallaxis per $206265''$ divisa accipienda est. Tandem nullo praecisionis detrimento in valoribus parallaxium pro $r$, $\lambda$, $\beta$ adhibere licebit $\Delta$, $l$, $b$, quoties in problemate inverso e loco parallaxi affecto locum ab eadem liberum determinare oportet.

\emph{Exemplum.} Sit ascensio recta Solis pro centro terrae $220^0\,46'\,44''{,}65=\lambda$, declinatio $-15^0\,49'\,43''{,}94=\beta$, distantia $0{,}9904311=r$; porro tempus sidereum in aliquo loco in terrae superficie gradibus expressum $78^0\,20'\,38''=L$, loci \SourceLeaf{091}elevatio poli $45^0\,27'\,57''=B$, parallaxis media solaris $8''{,}6=R$. Quaeritur locus Solis ex hoc loco visus, distantiaque ab eodem.
\begin{center}
\begin{tabular}{@{}ll@{}}
\makebox[15em][l]{$\log R$\hfill .\,.\,.\,.\,.\,.\,.\,.\,.} & $0{,}93450$\\[.3em]
\makebox[15em][l]{$\log\cos B$\hfill .\,.\,.\,.\,.\,.} & $9{,}84593$\\[.3em]
\makebox[15em][l]{$\mathrm{C.}\log r$\hfill .\,.\,.\,.\,.\,.\,.\,.} & $0{,}00418$\\[.3em]
\makebox[15em][l]{$\mathrm{C.}\log\cos\beta$\hfill .\,.\,.\,.\,.} & $0{,}01679$\\[.3em]
\makebox[15em][l]{$\log\sin(\lambda-L)$\hfill .\,.\,.} & $9{,}78509$\\[.3em]
\hline
\makebox[15em][l]{$\log(l-\lambda)$\hfill .\,.\,.\,.\,.\,.} & $0{,}58649$\\[.3em]
$l-\lambda=$ & $+3''{,}86$\\[.3em]
$l=$ & $220^0\,46'\,48''{,}51$\\[.3em]
\end{tabular}
\end{center}
\begin{center}
\begin{tabular}{@{}ll@{}}
\makebox[15em][l]{$\log\operatorname{tang}B$\hfill .\,.\,.\,.\,.} & $0{,}00706$\\[.3em]
\makebox[15em][l]{$\log\cos(\lambda-L)$\hfill .\,.} & $9{,}89909_n$\\[.3em]
\hline
\makebox[15em][l]{$\log\operatorname{tang}\theta$\hfill .\,.\,.\,.\,.\,.} & $0{,}10797_n$\\[.3em]
$\theta=$ & $127^0\,57'\quad0''$\\[.3em]
$\beta-\theta=$ & $-143^0\,46'\,44''$\\[.3em]
\end{tabular}
\end{center}
\begin{center}
\begin{tabular}{@{}ll@{}}
\makebox[15em][l]{$\log R$\hfill .\,.\,.\,.\,.\,.\,.\,.\,.} & $0{,}93450$\\[.3em]
\makebox[15em][l]{$\log\sin B$\hfill .\,.\,.\,.\,.\,.} & $9{,}85299$\\[.3em]
\makebox[15em][l]{$\mathrm{C.}\log r$\hfill .\,.\,.\,.\,.\,.\,.\,.} & $0{,}00418$\\[.3em]
\makebox[15em][l]{$\mathrm{C.}\log\sin\theta$\hfill .\,.\,.\,.\,.} & $0{,}10317$\\[.3em]
\makebox[15em][l]{$\log\sin(\beta-\theta)$\hfill .\,.\,.} & $9{,}77152_n$\\[.3em]
\hline
\makebox[15em][l]{$\log(b-\beta)$\hfill .\,.\,.\,.\,.\,.} & $0{,}66636_n$\\[.3em]
$b-\beta=$ & $-4''{,}64$\\[.3em]
$b=$ & $-15^0\,49'\,48''{,}58$\\[.3em]
\end{tabular}
\end{center}
\begin{center}
\begin{tabular}{@{}ll@{}}
\makebox[15em][l]{$\log(b-\beta)$\hfill .\,.\,.\,.\,.\,.} & $0{,}66636_n$\\[.3em]
$\log\operatorname{cotang}(\beta-\theta)$ & $0{,}13522$\\[.3em]
\makebox[15em][l]{$\log r$\hfill .\,.\,.\,.\,.\,.\,.\,.\,.\,.} & $9{,}99582$\\[.3em]
\makebox[15em][l]{$\log1''$\hfill .\,.\,.\,.\,.\,.\,.\,.\,.} & $4{,}68557$\\[.3em]
\hline
\makebox[15em][l]{$\log(r-\Delta)$\hfill .\,.\,.\,.\,.} & $5{,}48297_n$\\[.3em]
$r-\Delta=$ & $-0{,}000\,0304$\\[.3em]
$\Delta=$ & $0{,}990\,4615$\\[.3em]
\end{tabular}
\end{center}
